% Generated draft: requires mathematical and source review.
\documentclass[11pt]{article}
\usepackage[margin=25mm]{geometry}
\usepackage{fontspec}
\setmainfont{Latin Modern Roman}
\usepackage{amsmath,amssymb,mathtools,unicode-math}
\setmathfont{Latin Modern Math}
\usepackage{enumitem,needspace,xurl,hyperref}
\setlength{\parindent}{0pt}
\setlength{\parskip}{5pt}
\setlength{\emergencystretch}{3em}
\newcommand{\N}{\mathbb{N}}
\newcommand{\Z}{\mathbb{Z}}
\newcommand{\Q}{\mathbb{Q}}
\newcommand{\R}{\mathbb{R}}
\newcommand{\C}{\mathbb{C}}
\newcommand{\F}{\mathbb{F}}
\newcommand{\A}{\mathbb{A}}
\newcommand{\PP}{\mathbb{P}}
\newcommand{\E}{\mathbb{E}}
\newcommand{\eps}{\varepsilon}
\newcommand{\dd}{\,\mathrm{d}}
\begin{document}
\section*{UnsolvedMath ID: 800392}
\textbf{Source problem code:} catalogue entry
\subsection{Definitions and mathematical statement}
The supplied excerpt does not specify the ground field, its characteristic, or the full range of degrees. The following definitions give a conditional geometric formulation: let \(k\) be an algebraically closed field, let \(N\geq 1\), and let \(d\geq 2\). These are explicit working parameters requiring confirmation from the source context; the quoted results below are not asserted for every characteristic.

Let \(\operatorname{End}_d(\PP_k^N)\) denote the parameter space of morphisms \(f:\PP_k^N\to\PP_k^N\) represented by \(N+1\) homogeneous forms of degree \(d\) having no common geometric zero, modulo simultaneous nonzero scalar multiplication. Write
\[
\mathcal M_d^N=\operatorname{End}_d(\PP_k^N)/\operatorname{PGL}_{N+1}(k)
\]
for the moduli space under projective linear conjugacy, with \([f]\) denoting the class of \(f\). Define
\[
\operatorname{Aut}_k(f)=\{\gamma\in\operatorname{PGL}_{N+1}(k):\gamma^{-1}\circ f\circ\gamma=f\}.
\]
Its isomorphism type depends only on \([f]\). Here rationality over \(k\) means birational equivalence over \(k\) to an affine space. The source's notation \(\operatorname{Aut}(f)\) is interpreted geometrically in this formulation; whether the intended question instead concerns automorphisms defined over a specified nonclosed field is unresolved.

The supplied Question 15 asks:
\begin{enumerate}
\item[(a)] For \(N\geq 2\), determine which of the spaces \(\mathcal M_d^N\) are rational varieties. In particular, are all, some, or none of them rational, under the intended ground-field and degree hypotheses?

\item[(b)] For each pair \((N,d)\), describe all groups occurring as \(\operatorname{Aut}_k(f)\) for \([f]\in\mathcal M_d^N(k)\), or at least obtain a good upper bound, in terms of \(N\) and \(d\), for
\[
U(N,d)=\sup_{[f]\in\mathcal M_d^N(k)}\#\operatorname{Aut}_k(f).
\]
The excerpt reports, under its unstated ambient hypotheses, that \(U(N,d)\) is finite [MR2741188], that
\[
U(1,d)\leq\max\{60,2d+2\}
\]
[MR2884382, Example 2.54], and that
\[
U(2,d)\leq 60d^6
\]
[arxiv1509.06670, Theorem 6.2]. These are supplied background assertions, not independently verified claims for arbitrary \(k\).

\item[(c)] Describe the subvarieties corresponding to maps with nontrivial automorphism group; in particular, describe the geometry of the locus
\[
\mathcal A_d^N=\{[f]\in\mathcal M_d^N(k):\operatorname{Aut}_k(f)\neq\{1\}\}
\]
and the groups occurring on it. The excerpt cites a detailed analysis for \(N=1\) [MR3709645] and partial results for \(N=d=2\) [arxiv1607.05772]. It gives the example that, under the intended hypotheses, \(\mathcal M_2^1\cong\mathbb A_k^2\), and \(\mathcal A_2^1\) is a cuspidal cubic curve. Along that curve the automorphism group is \(\Z/2\Z\), except at the cusp, where it is \(\mathcal S_3\), the symmetric group on three letters. The characteristic restrictions needed for this example are not supplied.

\item[(d)] Describe the singular locus \(\operatorname{Sing}(\mathcal M_d^N)\), meaning the locus of points with nonregular local rings, for each intended pair \((N,d)\). The excerpt qualitatively states that most of these moduli spaces are singular and cites [MR3709645] for the case \(N=1\). It does not specify a quantified meaning of ``most.''
\end{enumerate}
\subsection{Short English statement}
Question 15 asks which moduli spaces of projective self-maps are rational, which automorphism groups occur and how large they can be, how maps with nontrivial automorphisms are distributed, and where the moduli spaces are singular. The ground field, characteristic, and full degree range need confirmation. The registry labels the problem open; this draft preserves that supplied label without independently assessing current status.
\subsection{Sources}
\begin{itemize}
\item Supplied local UnsolvedMath registry record, permanent ID \(800392\), exact source code \texttt{catalogue entry}, titled ``Open Problems in Arithmetic Dynamics --- Question 15,'' in the AMR Open Problem Lists collection. The registry attributes the question to Benedetto, DeMarco, Ingram, Jones et al., with proposed year 2018 and status \texttt{open}. Its creation and update timestamps are both \texttt{2026-10-16T14:00:00Z}; its supplied upstream revision is \texttt{d371ee917ac9647fbb06cd6fc529deec6f3dda82}.
\item Benedetto, DeMarco, Ingram, Jones et al., \emph{Current Trends and Open Problems in Arithmetic Dynamics} (2018), Question 15: \url{https://arxiv.org/abs/1806.04980}. The supplied provenance records access on 2026-10-16 and extraction by \texttt{arxiv\_latex}. This is supplied provenance, not a claim of independent access or verification.
\item [MR2741188], cited in the excerpt for finiteness of \(U(N,d)\). Full bibliographic details and applicable hypotheses are not supplied.
\item [MR2884382, Example 2.54], cited for \(U(1,d)\leq\max\{60,2d+2\}\). Full bibliographic details and applicable hypotheses are not supplied.
\item [arxiv1509.06670, Theorem 6.2], cited for \(U(2,d)\leq60d^6\). The supplied identifier is retained; no full bibliographic entry is supplied.
\item [MR3709645], cited for the analysis of automorphism loci and singularities when \(N=1\). Full bibliographic details are not supplied.
\item [arxiv1607.05772], cited for partial results on automorphism loci when \(N=d=2\). The supplied identifier is retained; no full bibliographic entry is supplied.
\end{itemize}
\end{document}
